% Préambule LaTeX de la version PDF (classe KOMA scrbook)

% --- légendes : les numéros de figures et tableaux sont ceux de l'original,
%     écrits dans le texte des légendes ; pas de numérotation automatique
\usepackage{caption}
\captionsetup{font=small, labelformat=empty, justification=justified,
  singlelinecheck=true, margin=0.6cm, skip=6pt}
\captionsetup[table]{position=top}

% --- titres
\addtokomafont{disposition}{\color{black!85}}
\addtokomafont{chapterprefix}{\mdseries\large\color{black!60}}
\addtokomafont{chapter}{\Huge}
\RedeclareSectionCommand[beforeskip=0pt, afterskip=2.2\baselineskip]{chapter}
\RedeclareSectionCommand[beforeskip=-2\baselineskip plus -1ex minus -.2ex,
  afterskip=0.8\baselineskip]{section}

% --- en-têtes et pieds de page
\usepackage[automark, headsepline=0.4pt]{scrlayer-scrpage}
\clearpairofpagestyles
\automark[chapter]{chapter}
\ohead{\headmark}
\ofoot*{\pagemark}
\setkomafont{pageheadfoot}{\normalfont\small\color{black!65}}
\setkomafont{pagenumber}{\normalfont\small}

% --- table des matières et liste des figures plus aérées
\RedeclareSectionCommand[tocindent=0pt]{chapter}


% --- typographie
\clubpenalty=10000
\widowpenalty=10000
\raggedbottom

% --- numérotation : seuls les chapitres sont numérotés par LaTeX ; les numéros
%     de paragraphes de l'original font partie des titres
\AtBeginDocument{\setcounter{secnumdepth}{0}}

% --- liste des figures : les numéros sont dans les légendes courtes
\captionsetup{listformat=empty}
\DeclareTOCStyleEntry[numwidth=0pt, indent=0pt, beforeskip=.15em]{tocline}{figure}
\DeclareTOCStyleEntry[numwidth=0pt, indent=0pt]{tocline}{table}
